% Study guide for the energy-methods / finite element series. Build: `mtk guide guide.tex` (after `mtk episodes`).
% Every number here is computed and asserted in check_guide.py (and the lesson's own numbers in solve.py).
\documentclass[11pt]{article}
\usepackage[margin=0.9in]{geometry}
\usepackage{amsmath,amssymb}
\usepackage{mtkguide}
\usepackage[hidelinks]{hyperref}
\setlength{\parindent}{0pt}
\setlength{\parskip}{0.4em}
\newcommand{\mm}{\ \mathrm{mm}}
\newcommand{\Nmm}{\ \mathrm{N/mm}}

\begin{document}
{\LARGE\bfseries Energy Methods and the Finite Element Method}\par{\large Practice Guide}\par
{\color{mtkmuted} AERSP 301, HW 6 topics \textperiodcentered{} a guide to work through with the videos}\par\medskip

\textbf{How to use it.} The video pauses at each checkpoint. Work it here, on paper or a tablet, before you press
play: the video then works the same step, so you find out at once whether you had it. Stuck for five minutes? Read
the hint on page~\pageref{hints}. After the last video, two fresh practice problems (page~\pageref{practice}) repeat
both homework methods with new numbers; their full solutions are at the end.

Notation, as in the videos: $U$ strain energy, $W$ work potential of the loads, $\Pi = U - W$ total potential
energy; $q_i$ node displacements; a spring's stretch is (right-end displacement) $-$ (left-end displacement).

\section*{Episode 1: Work, Energy and the Bowl}

\checkpoint[3cm]{1.1}{The bottom of the bowl}{%
A spring of stiffness $K = 20\Nmm$ is fixed to a wall and its free end carries a load $F = 100\ \mathrm{N}$. Its
total potential energy is $\Pi(q) = \tfrac12 K q^2 - F q$. Write $\Pi(q)$ with the numbers in, then find the $q$ that
makes $\Pi$ smallest, and the value of $\Pi$ there.}

\section*{Episode 2: Many Springs}

\checkpoint[3.5cm]{2.1}{Two springs in a row}{%
Wall $-$ spring $k_1 = 40\Nmm$ $-$ node 1 $-$ spring $k_2 = 20\Nmm$ $-$ node 2. Loads $F_1 = 40\ \mathrm{N}$ at node 1
and $F_2 = 20\ \mathrm{N}$ at node 2, both to the right. Solve
$\begin{bmatrix}60&-20\\-20&20\end{bmatrix}\begin{Bmatrix}q_1\\q_2\end{Bmatrix}=\begin{Bmatrix}40\\20\end{Bmatrix}$
for $q_1$ and $q_2$ (mm).}

\checkpoint[3cm]{2.2}{Row 3 of the four-node matrix}{%
Node 1 (on the wall) $-$ $k_1$ $-$ node 2 $-$ two springs $k_2$ side by side $-$ node 3 $-$ $k_3$ $-$ node 4. Using the
rule ``a spring $k$ between nodes $i$ and $j$ adds $+k$ at $(i,i)$ and $(j,j)$ and $-k$ at $(i,j)$ and $(j,i)$'', write
row 3 of the $4\times4$ stiffness matrix.}

\section*{Episode 3: Homework Problem 1}

\checkpoint[4cm]{3.1}{The derivative for the bar}{%
For homework problem 1,
$\Pi = \tfrac12 K_1 q_2^2 + \tfrac12 K_2 q_1^2 + \tfrac12 K_2 (q_2-q_1)^2 + \tfrac12 K_3 (q_3-q_2)^2
- P_1 q_1 - (P_2+P_3) q_2 - P_4 q_3$.
Find $\partial \Pi / \partial q_2$, set it to zero, and collect the terms by $q_1$, $q_2$, $q_3$.}

\checkpoint[3cm]{3.2}{Test the answer with numbers}{%
The video found $q_2 = \dfrac{P_1 + 2(P_2+P_3+P_4)}{2K_1+K_2}$, $q_1 = \dfrac{P_1 + K_2 q_2}{2K_2}$ and
$q_3 = q_2 + \dfrac{P_4}{K_3}$. Put in $K_1 = K_2 = K_3 = 100\Nmm$ and $P_1 = P_2 = P_3 = P_4 = 30\ \mathrm{N}$.
Then check one spring with Newton: does $K_3$ carry exactly $P_4$?}

\section*{Episode 4: From Springs to Bars}

\checkpoint[3cm]{4.1}{How good is the straight line?}{%
A bar fixed at $x = 0$ carries a uniform load $f$ per unit length. The assumed shape $u = a x$ gave
$a = \dfrac{fL}{2EA}$. The exact answer is $u(x) = \dfrac{f}{EA}\left(Lx - \dfrac{x^2}{2}\right)$. Compare the two at
the middle, $x = L/2$, in units of $fL^2/EA$.}

\section*{Episode 5: One Finite Element}

\checkpoint[2cm]{5.1}{An element's stiffness}{%
An aluminum element has $E = 70\,000\ \mathrm{N/mm^2}$, $A = 100\mm^2$ and length $l = 250\mm$. What is its stiffness
$EA/l$?}

\section*{Episode 6: Assembly and Solving}

\checkpoint[3.5cm]{6.1}{The column's load vector}{%
A column of height $L$ stands on the ground and carries its own weight, $c = -\rho A g$ per unit length. It is split
into three equal elements ($l = L/3$), nodes 1 (base) to 4 (top). Each element sends $cl/2$ to each of its nodes, and
the ground pushes on node 1 with a reaction $R$. Write the global load vector $\{F_1, F_2, F_3, F_4\}$.}

\checkpoint[3cm]{6.2}{Stress in element 2}{%
The solution was $\{q_2, q_3, q_4\} = \dfrac{\rho g L^2}{2E}\{-5/9,\,-8/9,\,-1\}$, with $q_1 = 0$. Using
$\sigma = E\,(q_{\text{top}} - q_{\text{bottom}})/l$ and $l = L/3$, find the stress in element 2 in units of
$\rho g L$.}

\section*{Episode 7: A Stepped Bar and a Spring}

\checkpoint[5cm]{7.1}{Assemble four by four}{%
Nodes 0 (left wall), 1 (the step), 2, 3 (right wall). Element 1 joins nodes 0 and 1 with $k_1 = 28\,000\Nmm$;
element 2 joins nodes 1 and 2 with $k_2 = 14\,000\Nmm$; a spring $K = 14\,000\Nmm$ joins nodes 2 and 3. Assemble the
$4\times4$ global stiffness matrix.}

\checkpoint[2.5cm]{7.2}{Stress in the thin piece}{%
With $E = 70\,000\ \mathrm{N/mm^2}$, $l = 500\mm$, $q_1 = 0.2\mm$ and $q_2 = 0.1\mm$, find the stress in element 2
(nodes 1 to 2). Tension or compression?}

\section*{Episode 8: Homework Problem 2}

\checkpoint[5cm]{8.1}{Homework 2: K and F}{%
Two elements of $l = 20$ in, each with $EA/l = 1000$ lb/in: node 0 (wall, $x = 0$), node 1 ($x = 20$ in,
$P_1 = 50$ lb to the right), node 2 ($x = 40$ in, $P_2 = 100$ lb to the right). A uniform load of 5 lb/in points to
the \emph{left} along the whole bar. Assemble the $3\times3$ global stiffness matrix $K$ and the load vector $F$
(include the wall reaction $R$ at node 0).}

\checkpoint[2.5cm]{8.2}{Homework 2: the reaction}{%
The solution is $q_0 = 0$, $q_1 = 0$, $q_2 = 0.05$ in. Use row 0 of $Kq = F$ (the constrained equation) to find the
wall reaction $R$. Which way does the wall push?}

\section*{After the video: two practice problems}\label{practice}

\textbf{Practice A (springs, like homework problem 1).} Wall $-$ spring $k_1 = 20\Nmm$ $-$ node 1 $-$ spring
$k_2 = 20\Nmm$ $-$ node 2 $-$ spring $k_3 = 10\Nmm$ $-$ a second wall. A load $F_2 = 60\ \mathrm{N}$ pulls node 2 to
the right; node 1 carries no load. (a) Write $U$, $W$ and $\Pi$. (b) Set $\partial\Pi/\partial q_1 = 0$ and
$\partial\Pi/\partial q_2 = 0$ and solve. (c) Find the force each wall exerts, and check that they balance the load.

\textbf{Practice B (finite elements, like homework problem 2).} A steel bar, $E = 200\,000\ \mathrm{N/mm^2}$,
$A = 60\mm^2$, $L = 1200\mm$, is fixed at $x = 0$. A uniform load $f = 1.5\Nmm$ acts along $+x$, and a tip load
$P = 300\ \mathrm{N}$ acts along $+x$. Use three equal elements. (a) Find each element's stiffness and load vector.
(b) Assemble $K$ and $F$. (c) Apply the support and solve for the three free node displacements. (d) Find the stress
in each element. (e) Use the constrained equation to find the wall reaction.

\newpage
\section*{Hints}\label{hints}
\hintfor{1.1} Differentiate $\Pi$ with respect to $q$ and set the slope to zero.
\hintfor{2.1} Add the two equations: one unknown cancels.
\hintfor{2.2} Row 3 belongs to node 3. Which springs touch node 3, and what is at the other end of each?
\hintfor{3.1} $q_2$ appears in three springs. The chain rule gives a factor $-1$ where $q_2$ is subtracted.
\hintfor{3.2} Work out $q_2$ first; the other two follow from it.
\hintfor{4.1} Evaluate both expressions at $x = L/2$.
\hintfor{5.1} Plug in; watch the units: $\mathrm{N/mm^2}\times\mathrm{mm^2}/\mathrm{mm} = \mathrm{N/mm}$.
\hintfor{6.1} A node shared by two elements receives a share from each.
\hintfor{6.2} Element 2 runs from node 2 to node 3.
\hintfor{7.1} Each spring or element drops a $2\times2$ block on its two nodes; where blocks meet on the diagonal,
add.
\hintfor{7.2} $\sigma = E\,(q_b - q_a)/l$ with $q_b$ the right node.
\hintfor{8.1} Each element sends $cl/2 = (-5)(20)/2$ to each of its nodes; add the point loads at their nodes.
\hintfor{8.2} Row 0 reads $1000\,q_0 - 1000\,q_1 = R - 50$.

\section*{Solutions}\label{solutions}

\solutionfor{1.1}{The bottom of the bowl}
$\Pi = 10q^2 - 100q$ (N mm). $d\Pi/dq = 20q - 100 = 0$, so $q = 5\mm$, and $\Pi(5) = 250 - 500 = -250$ N mm.
Check: $Kq = 20\times 5 = 100 = F$, Newton's balance.

\solutionfor{2.1}{Two springs in a row}
Adding the rows: $40q_1 = 60$, so $q_1 = 1.5\mm$. Row 2: $20(q_2 - 1.5) = 20$, so $q_2 = 2.5\mm$.
Check: spring 1 carries $40 + 20 = 60$ N and stretches $60/40 = 1.5$; spring 2 carries 20 N and stretches
$20/20 = 1 = q_2 - q_1$.

\solutionfor{2.2}{Row 3 of the four-node matrix}
Node 3 touches the pair of $k_2$ springs (to node 2) and $k_3$ (to node 4). Row 3 $= [\,0,\ -2k_2,\ 2k_2 + k_3,\ -k_3\,]$.
It adds up to zero, as every row does before supports are removed.

\solutionfor{3.1}{The derivative for the bar}
$\dfrac{\partial\Pi}{\partial q_2} = K_1 q_2 + K_2(q_2 - q_1) - K_3(q_3 - q_2) - (P_2 + P_3) = 0$, so
$-K_2 q_1 + (K_1 + K_2 + K_3)\,q_2 - K_3 q_3 = P_2 + P_3$.

\solutionfor{3.2}{Test the answer with numbers}
$q_2 = (30 + 2\cdot 90)/300 = 0.7\mm$; $q_1 = (30 + 70)/200 = 0.5\mm$; $q_3 = 0.7 + 30/100 = 1.0\mm$.
$K_3$ stretches $q_3 - q_2 = 0.3\mm$ and carries $100\times0.3 = 30$ N $= P_4$. The wall springs carry
$K_1 q_2 + K_2 q_1 = 70 + 50 = 120$ N $= 4\times30$: everything balances.

\solutionfor{4.1}{How good is the straight line?}
Assumed: $u(L/2) = aL/2 = fL^2/(4EA) = 0.25\,fL^2/EA$. Exact: $\dfrac{f}{EA}\left(\dfrac{L^2}{2} - \dfrac{L^2}{8}\right)
= 0.375\,fL^2/EA$. The straight line is too low inside the bar, though it matches at the tip.

\solutionfor{5.1}{An element's stiffness}
$EA/l = 70\,000\times100/250 = 28\,000\Nmm$.

\solutionfor{6.1}{The column's load vector}
$F = \dfrac{cl}{2}\{1,\,2,\,2,\,1\} + \{R,\,0,\,0,\,0\}$: the shared middle nodes get $cl/2$ from each side.

\solutionfor{6.2}{Stress in element 2}
$q_3 - q_2 = (-8/9 + 5/9)\,\rho g L^2/(2E) = -\tfrac13\,\rho g L^2/(2E)$. Times $E/l = 3E/L$:
$\sigma_2 = -\tfrac12\,\rho g L$ (compression). It equals the exact stress $-\rho g L(1 - x/L)$ at the element's middle,
$x = L/2$.

\solutionfor{7.1}{Assemble four by four}
$K = \begin{bmatrix}28\,000&-28\,000&0&0\\-28\,000&42\,000&-14\,000&0\\0&-14\,000&28\,000&-14\,000\\0&0&-14\,000&14\,000\end{bmatrix}$ N/mm.
The overlaps add: $28\,000 + 14\,000$ at node 1 and $14\,000 + 14\,000$ at node 2.

\solutionfor{7.2}{Stress in the thin piece}
$\sigma_2 = 70\,000\,(0.1 - 0.2)/500 = -14\ \mathrm{N/mm^2} = -14$ MPa: compression.

\solutionfor{8.1}{Homework 2: K and F}
$K = 1000\begin{bmatrix}1&-1&0\\-1&2&-1\\0&-1&1\end{bmatrix}$ lb/in. Each element sends $(-5)(20)/2 = -50$ lb to each
node, so $F = \{R - 50,\ -50 - 50 + 50,\ -50 + 100\} = \{R - 50,\ -50,\ 50\}$ lb.

\solutionfor{8.2}{Homework 2: the reaction}
$1000(0) - 1000(0) = R - 50$, so $R = 50$ lb, positive: the wall pushes the bar in $+x$.
Check: $50 + 50 + 100 - 5\times40 = 0$.

\subsection*{Practice A}
(a) $U = \tfrac12(20)q_1^2 + \tfrac12(20)(q_2 - q_1)^2 + \tfrac12(10)q_2^2$ (the right spring's stretch is $0 - q_2$,
and squared it is $q_2^2$); $W = 60\,q_2$; $\Pi = U - W$.
(b) $\partial\Pi/\partial q_1 = 20q_1 - 20(q_2 - q_1) = 0$ and $\partial\Pi/\partial q_2 = 20(q_2 - q_1) + 10q_2 - 60 = 0$,
i.e. $\begin{bmatrix}40&-20\\-20&30\end{bmatrix}\begin{Bmatrix}q_1\\q_2\end{Bmatrix} = \begin{Bmatrix}0\\60\end{Bmatrix}$.
The first row gives $q_2 = 2q_1$; then $-20q_1 + 60q_1 = 60$, so $q_1 = 1.5\mm$, $q_2 = 3.0\mm$.
(c) Left wall: $-k_1 q_1 = -30$ N (it pulls the stretched spring back, to the left). Right wall: the right spring is
squeezed by $3.0\mm$, $10\times3 = 30$ N, and the wall pushes left. $60 - 30 - 30 = 0$. Also $\Pi = -90$ N mm
$= -\tfrac12(60)(3.0)$.

\subsection*{Practice B}
(a) $l = 400\mm$, $k = EA/l = 200\,000\times60/400 = 30\,000\Nmm$; each element's load vector is
$cl/2 = 1.5\times400/2 = 300$ N at each node.
(b) $K = 30\,000\begin{bmatrix}1&-1&0&0\\-1&2&-1&0\\0&-1&2&-1\\0&0&-1&1\end{bmatrix}$,
$F = \{300 + R,\ 600,\ 600,\ 300 + 300\} = \{300 + R,\ 600,\ 600,\ 600\}$ N.
(c) Remove row and column 0. Last row: $30\,000(q_3 - q_2) = 600$, so $q_3 - q_2 = 0.02$. Adding it to row 2:
$30\,000(q_2 - q_1) = 1200$, so $q_2 - q_1 = 0.04$. Adding both to row 1: $30\,000\,q_1 = 1800$, so $q_1 = 0.06\mm$;
then $q_2 = 0.10\mm$, $q_3 = 0.12\mm$. Exact tip: $PL/EA + fL^2/(2EA) = 0.03 + 0.09 = 0.12\mm$.
(d) $\sigma = E(q_b - q_a)/l$: $200\,000\times0.06/400 = 30$, then $20$, then $10$ MPa, all tension.
(e) Row 0: $30\,000(0 - 0.06) = 300 + R$, so $R = -2100$ N: the wall pulls left with the whole load,
$fL + P = 1800 + 300$.

\end{document}
